\chapter{Problem Definition and Mathematical Formulation}
\label{chap:problem}

% ---------------------------------------------------------------
\section{Chapter Overview and Notation}
% ---------------------------------------------------------------
\begin{table}[htbp]
\centering
\small
\ifshowrevisions\color{P3Rcolor}\fi
\begin{tabular}{l>{\raggedright\arraybackslash}p{9.8cm}}
\toprule
\textbf{Symbol} & \textbf{Description} \\
\midrule
$t$ & Hourly timestep ($\Delta t=1\,\mathrm{h}$) \\
$P_t^{\PV},P_t^{\Load}$ & PV generation and telecom load (kW) \\
$G_t\in\{0,1\}$ & Grid availability from the site trace in the primary campaign \\
$P_t^{\Grid}$ & Total grid power, including any battery charging (kW) \\
$P_t^{\Grid,\max}$ & Maximum grid connection capacity (kW) \\
$\SoC_t$ & Physical battery state of charge in $[0.20,1.00]$; normalised for the agent \\
$E_{\Bat}$ & Battery energy capacity (kWh) \\
$P_t^{\Bat}$ & Net battery power (positive: charging; negative: discharging), kW \\
$P_{\max}^{\ch},P_{\max}^{\dis}$ & Battery charge/discharge power limits (kW) \\
$\eta_{\ch},\eta_{\dis}$ & Site-data battery charge/discharge coefficients used by the simulator \\
$\DG_t\in\{0,1\}$ & DG ON/OFF request \\
$P_{\mathrm{rated}}^{\DG}$ & DG rated power limit (kW) \\
$E_t^{\DG}$ & Actual diesel-generated energy delivered to the residual load (kWh/step) \\
$\Inv_t$ & On-site inventory in kWh-equivalent units \\
$\Pipe_t$ & Pending diesel quantity in kWh-equivalent units; at most one order \\
$F_t^{\mathrm{pipe}}$ & Binary pending-order flag \\
$Q_t^{\mathrm{pipe}}$ & Pending-order quantity (kWh-equivalent) \\
$Q_t$ & New order quantity (kWh-equivalent) \\
$C^{\Tank}$ & Tank capacity ($72\bar D$ kWh-equivalent in the baseline) \\
$r_{\DG}$ & Reporting conversion factor, $0.30$ litres/kWh; not used in inventory transitions \\
$A_t$ & Diesel arrival added to inventory (kWh-equivalent) \\
$L_t$ & Delivery lead time (hourly steps) \\
$\tau_t$ & Hours elapsed since the last order was placed \\
$H_t^{\mathrm{rem}}$ & Delivery-remaining observation, zero in the primary campaign \\
$\hour_t$ & Hour-of-day encoded by sine and cosine (two scalars) \\
$U_t$ & Unserved electrical energy (kWh/step) \\
$\bar D$ & Mean hourly site load of the data slice supplied to the simulator (kWh) \\
$S_t,V_t$ & Stockout and physical-violation indicator variables \\
$\alpha,\beta,\lambda$ & Diesel, grid and unserved-energy reward weights \\
$\gamma_r,\mu$ & Stockout and physical-violation penalties \\
$\gamma$ & RL discount factor, implemented value $0.99$ \\
\bottomrule
\end{tabular}
\caption{Notation for the implemented formulation and control model.}
\label{tab:notation}
\end{table}

As established in Chapter~\ref{chap:literature} and summarised in
Table~\ref{tab:litgap}, to the best of the author's knowledge no prior work
jointly models diesel fuel as a finite stochastic inventory while
simultaneously optimising operational energy dispatch in telecom towers.
In rural deployments, diesel logistics---including pilferage, stockouts, and
delivery disruptions---contribute significantly to tower outages
\citep{datoms2024,alton2025,gsma2025b,trai2024}.

This chapter formulates the coupled diesel-inventory and tower-energy
control problem as a feasibility-constrained sequential decision process.
The learned actions are generator commitment and diesel replenishment;
PV, grid and battery dispatch follow the simulator's fixed priority rule.
Decisions are hourly ($\Delta t=1\,\mathrm{h}$), consistent with the
ITU/Zindi site traces described in Chapter~\ref{chap:datasets-setup}.


% ---------------------------------------------------------------
\section{System Model}
% ---------------------------------------------------------------

The overall hybrid telecom tower and diesel-logistics system is illustrated in
Figure~\ref{fig:system_diagram}.

\begin{figure}[htbp]
\centering
\resizebox{0.9\textwidth}{!}{%
\begin{tikzpicture}[
    >=Stealth,
    node distance=2.6cm and 3.2cm,
    box/.style={draw, rounded corners=4pt, minimum height=1.1cm,
                minimum width=2.8cm, align=center, font=\small},
    agent/.style={draw, fill=gray!12, rounded corners=6pt,
                  minimum width=4.5cm, minimum height=1.4cm,
                  align=center, font=\bfseries\small},
    arrow/.style={->, thick},
    obs/.style={->, dashed, thick, blue!70!black},
    act/.style={->, dashed, thick, red!70!black}
]

% --- Energy components (top row) ---
\node[box] (pv)   {Solar PV};
\node[box, right=of pv]   (grid) {Grid Supply\\(intermittent)};
\node[box, right=of grid] (bat)  {Battery Storage\\(LFP)};
\node[box, right=of bat]  (dg)   {Diesel Generator\\(\revised{rated $P_{\mathrm{rated}}^{\DG}$})};
\node[box, right=of dg]   (load) {Telecom Load\\(BTS + ancillary)};

% --- RL agent ---
\node[agent, below=2.8cm of bat] (agent) {RL Agent};

% --- Diesel logistics (bottom row) ---
\node[box, below=5.0cm of dg] (tank)
  {Diesel Tank\\(\revised{$\Inv_t$ kWh-equiv.})};
\node[box, left=of tank]       (pipe) {Pipeline\\($\Pipe_t$)};
\node[box, left=of pipe]       (sup)  {Fuel Supplier};

% --- Physical flows ---
\draw[arrow] (pv.south east)
  -- ++(0,-0.55cm) -| (bat.south west);
\draw[arrow] (grid) -- (bat);
\draw[arrow] (bat.north)
    -- ++(0,0.45cm) -| (load.north east);
% Direct PV/grid supply to the load (in addition to battery charging).
\ifshowrevisions
  \draw[arrow,draw=P3Rcolor] (pv.north) -- ++(0,1.0) -| (load.north west);
  \draw[arrow,draw=P3Rcolor] (grid.north) -- ++(0,1.45) -| (load.north);
\else
  \draw[arrow] (pv.north) -- ++(0,1.0) -| (load.north west);
  \draw[arrow] (grid.north) -- ++(0,1.45) -| (load.north);
\fi
\draw[arrow] (dg)   -- (load);
\draw[arrow] (sup)  -- node[above,font=\footnotesize] {$Q_t$ (order)} (pipe);
\draw[arrow] (pipe) -- node[above,font=\footnotesize] {$A_t$ (arrival)} (tank);
\draw[arrow] (tank) -- node[right,font=\footnotesize] {fuel} (dg);


% --- Observations ---
\draw[blue, dashed, ->] (pv.south east) to[out=-70,in=160] (agent.west);
\draw[blue, dashed, ->] (grid.south) to[out=-90,in=120] (agent.north west);
\draw[blue, dashed, ->] (bat.south west) to[out=-120,in=60] (agent.north east);
\draw[blue, dashed, ->] (load.south west) to[out=-130,in=20] (agent.east);
\draw[blue, dashed, ->] (pipe.north) to[out=90,in=-90] (agent.south);
\draw[obs] (tank.north west) -- (agent.south east);

\node[
    font=\scriptsize,
    above=7mm of agent,
    align=left,
    fill=white,
    inner sep=1pt
]{
  \textbf{State $s_t$ (selected):} PV output $P_t^{\PV}$, load $P_t^{\Load}$,\\
  grid availability $G_t$, battery SoC,\\
  diesel inventory $\Inv_t$, pipeline $\Pipe_t$,\\
  time-of-day encoding $\hour_t$
};

% --- Actions ---
\draw[act] (agent.60)  -- node[right,font=\footnotesize] {$\DG_t \in \{0,1\}$}    (dg.south);
\draw[act] (agent.240) -- node[left,font=\footnotesize]  {$Q_t$}                   (sup.north);

\end{tikzpicture}%
}
\caption{System architecture of the hybrid telecom tower under DG ON/OFF
modelling. The RL agent observes energy states (PV, load, grid, battery SoC)
and diesel logistics states (tank inventory $\Inv_t$ and in-transit pipeline
$\Pipe_t$), and requests DG commitment and tank-relative replenishment.
Actual DG energy supplied is limited by rated capacity and residual
load. PV and grid can supply the load and charge the battery; battery
dispatch is not an RL action.}
\label{fig:system_diagram}
\end{figure}

The system comprises rooftop solar PV, a lithium iron phosphate (LFP) battery
bank, intermittent grid supply, a diesel generator \revised{subject to a fixed rated-power limit}
$P_{\mathrm{rated}}^{\DG}$ when switched on, and a finite diesel tank
replenished by external truck deliveries subject to stochastic lead times.
$G_t \in \{0,1\}$ is a binary grid availability indicator; when $G_t = 0$ the
grid is unavailable due to outage and no power can be drawn from it.
Decisions are taken every hour ($\Delta t = 1$\,h).
\begin{revisedblock}
The operating assumptions are non-cancellable replenishment orders,
at most one pending order, and a generator fuel-start reserve equal to
$1.20\times2\,\mathrm{h}\times P_{\mathrm{rated}}^{\DG}$.
The reserve is a feasibility threshold, not a two-hour minimum
continuous ON-time. Future delivery times and supplier regimes are
not observed by the controller.
The study models ten towers independently; multi-tower coordination,
shared supplier capacity and field deployment are outside its scope.
\end{revisedblock}

At each timestep, \revised{the power balance accounts for
the energy supplied to the load and any unmet demand:}
\begin{equation}
P_t^{\Load} =
\revised{
P_t^{\PV,\mathrm{used}}
- P_t^{\Bat}
+ P_t^{\Grid}
+ P_t^{\DG,\mathrm{used}}
+ \frac{U_t}{\Delta t}
},
\label{eq:balance}
\end{equation}
\revised{where $P_t^{\PV,\mathrm{used}}$ denotes PV power
used by the load or battery, $P_t^{\DG,\mathrm{used}}$
denotes actual DG power supplied, and $U_t$ is unmet
load energy (kWh). Grid draw satisfies
$0 \leq P_t^{\Grid} \leq G_t P_t^{\Grid,\max}$,
while DG supply is bounded by
$0 \leq P_t^{\DG,\mathrm{used}}
\leq \DG_t P_{\mathrm{rated}}^{\DG}$.}
The negative sign on $P_t^{\Bat}$ reflects the convention
of Table~\ref{tab:notation}: a positive value denotes
charging (power absorbed from the bus), so a discharging
battery contributes positively to the supply side via
the negative sign.

Although diesel inventory decisions operate on multi-day timescales, they
directly constrain hourly dispatch feasibility. During grid outages or
low-solar periods, the generator can only operate if sufficient fuel remains
on site. A controller that optimises hourly energy dispatch without accounting
for inventory state may select myopic dispatch decisions that later lead to infeasible inventory states. Jointly modelling inventory and dispatch
decisions creates the opportunity for the agent to anticipate future fuel constraints and maintain
service reliability under uncertain delivery schedules.

The formulation therefore differs from conventional telecom energy-management
models in two respects. First, diesel fuel is represented as a stochastic
inventory subject to delivery delays rather than as an unconstrained fuel
source. Second, dispatch and inventory decisions are optimised jointly within
a single sequential decision framework, allowing the controller to anticipate
future fuel availability when making present operational decisions.

% ---------------------------------------------------------------
\section{Motivation for Reinforcement Learning}
\label{sec:rl_motivation}
% ---------------------------------------------------------------

\begin{revisedblock}
The ordering decision influences generator feasibility across many
hourly steps, while delivery times are stochastic. The primary
experiments use recorded PV, load and grid traces with simulated
supplier lead times. Reinforcement learning provides one approach
to selecting generator and replenishment actions without solving a
fresh multi-period optimisation problem at every step. The implemented
policy and its alternatives are detailed in Chapter~\ref{chap:methodology}.
\end{revisedblock}

% ---------------------------------------------------------------
\section{State Space}
\label{sec:state}
% ---------------------------------------------------------------

\begin{revisedblock}
The controller receives an 11-scalar observation:
\begin{equation}
\begin{aligned}
s_t = [\,&\SoC_t,\ \Inv_t,\ F_t^{\mathrm{pipe}},\ Q_t^{\mathrm{pipe}},
\ P_t^{\PV},\ P_t^{\Load},\ G_t,\\
&\sin(2\pi h_t/24),\ \cos(2\pi h_t/24),\
\tau_t,\ H_t^{\mathrm{rem}}\,].
\end{aligned}
\label{eq:state}
\end{equation}
Here $F_t^{\mathrm{pipe}}$ indicates whether a delivery is pending,
$Q_t^{\mathrm{pipe}}$ is its quantity, and $\tau_t$ counts the hours
since the most recent order. The scalar channels are normalised
where applicable before policy evaluation. The remaining-delivery
channel $H_t^{\mathrm{rem}}$ is zero in the principal experiments;
an exploratory EWMA-based supplier-delay experiment uses an additional
twelfth channel, as described in Chapter~\ref{chap:results}.
\end{revisedblock}

\begin{table}[htbp]
\centering
\small
\ifshowrevisions\color{P3Rcolor}\fi
\begin{tabularx}{\textwidth}{p{2.3cm}p{3.0cm}X}
\toprule
\textbf{Channel} & \textbf{Source} & \textbf{Notes} \\
\midrule
$P_t^{\PV}$ & ITU/Zindi & Hourly solar trace \\
$P_t^{\Load}$ & ITU/Zindi & Hourly demand trace \\
$G_t$ & ITU/Zindi & Recorded availability in primary evaluation \\
$\SoC_t$ & Simulator & Initialised from site parameters; updated by fixed dispatch \\
$\Inv_t$ & Simulated & Initialised as a fraction of assumed tank capacity; not measured fuel-level data \\
$F_t^{\mathrm{pipe}}$ & Simulator & Pending-order indicator \\
$Q_t^{\mathrm{pipe}}$ & Simulator & Outstanding order quantity \\
$\tau_t$ & Simulator & Elapsed hours since order placement \\
$H_t^{\mathrm{rem}}$ & Inactive channel & Zero in principal experiments \\
$\sin,\cos$ hour & Derived & From site timestamp, two scalar channels \\
\bottomrule
\end{tabularx}
\caption{Data provenance of the 11 observation channels. Inventory,
order status, delivery events and ordering history are simulated;
they are not measured fuel logistics in the source dataset.}
\label{tab:state_sources}
\end{table}

The pipeline quantity is capped at one pending order at any time: if a
delivery is already in transit, the feasible action set restricts $Q_t = 0$
until the pending order arrives (Section~\ref{sec:constraints}).

% ---------------------------------------------------------------
\section{Action Space}
\label{sec:action}
% ---------------------------------------------------------------

\begin{revisedblock}
At each hourly step the controller selects a generator request
$\DG_t\in\{0,1\}$ and an order class $k_t\in\{0,1,2\}$.
The joint action is
\begin{equation}
a_t=(\DG_t,k_t),\qquad
\operatorname{index}(a_t)=3\DG_t+k_t,
\label{eq:action}
\end{equation}
with order quantity
\[
Q_t=\begin{cases}
0,&k_t=0,\\
0.30\,C^{\Tank},&k_t=1,\\
0.60\,C^{\Tank},&k_t=2.
\end{cases}
\]
Thus the environment has six discrete joint actions
(\texttt{Discrete(6)}), rather than fixed-litre tanker sizes.
A pending order blocks further ordering; insufficient tank
headroom restricts or clips the requested quantity according to
the action-feasibility interface. DG activation is subject to the
fuel-start reserve. Battery dispatch follows the fixed environment
rule, not the learned action.
\end{revisedblock}

% ---------------------------------------------------------------
\section{System Dynamics and Uncertainty}
\label{sec:dynamics}
% ---------------------------------------------------------------

\paragraph{Battery dynamics.}
Battery state-of-charge evolves according to
\begin{equation}
\SoC_{t+1} = \SoC_t
  + \frac{\eta_{\ch}\,\max(P_t^{\Bat},0)\,\Delta t
          - \max(-P_t^{\Bat},0)\,\Delta t\,/\,\eta_{\dis}}
         {E_{\Bat}}.
\label{eq:soc}
\end{equation}
\revised{The coefficients $\eta_{\ch}$ and $\eta_{\dis}$ are taken
from the site data and used by the simulator as implemented;
they should not be assumed to be conventional physical efficiencies.
The site-data coefficient convention is discussed in
Chapter~\ref{chap:methodology}.}
\paragraph{Diesel inventory dynamics.}
On-site fuel inventory evolves as
\begin{revisedblock}
\begin{equation}
\Inv_{t+1}
=\min\!\left(C^{\Tank},\,
\max(0,\Inv_t-E_t^{\DG})+A_t\right),
\label{eq:inv}
\end{equation}
\end{revisedblock}
where the arrival term is
\begin{revisedblock}
\begin{equation}
A_t=\begin{cases}
Q_t^{\mathrm{pipe}}, & \text{if the outstanding order arrives at step }t,\\
0, & \text{otherwise,}
\end{cases}
\label{eq:arrival}
\end{equation}
\end{revisedblock}

\paragraph{Pipeline dynamics.}
\begin{revisedblock}
The pipeline state records the quantity in transit, represented in
the observation by a pending flag and pending quantity. At most one
order may be outstanding. Each step consumes available DG energy
before processing arrival of an existing pending order; a new order
is accepted afterwards and cannot arrive during that same step.
\end{revisedblock}

\paragraph{Stochastic lead time.}
In the primary experimental campaign (Chapter~\ref{chap:results}), delivery
is modelled as a memoryless geometric process: once an order is placed,
in each subsequent hour that it has not yet arrived, delivery occurs
independently with probability $p$, so that
\begin{equation}
L_t \;\sim\; \Geometric(p), \qquad \mathbb{E}[L_t] = 1/p \text{ hours},
\label{eq:leadtime}
\end{equation}
consistent with $\Delta t = 1$\,h. The four evaluated logistics scenarios
correspond to different mean lead times, calibrated through $p$: \emph{Normal}
($p=0.04167$, mean 24\,h), \emph{Delayed} ($p=0.02083$, mean 48\,h),
\emph{Monsoon} ($p=0.01389$, mean 72\,h), and \emph{Extreme} ($p=0.00298$,
mean 336\,h, $\approx$14 days). All four scenarios are evaluated in
Chapter~\ref{chap:results} to test robustness under increasing
supply-chain disruption.

As a separate robustness check reported alongside the primary results
(Chapter~\ref{chap:results}), rather than as part of the primary
comparison, delivery lead time is also modelled using log-normal and
Weibull distributions with the same mean lead time as the corresponding
geometric setting, using shape parameters $\sigma=0.5$ (moderate
variability) and $\sigma=0.8$ (heavy-tailed) for the log-normal case, and
$k=2$ for the Weibull case. This tests whether the qualitative policy
ranking is sensitive to the specific parametric form assumed for delivery
uncertainty, independent of the primary geometric model used elsewhere in
this thesis. The log-normal distribution is standard in remote diesel
logistics modelling \citep{srinivasan2013,shen2011}.

\begin{revisedblock}
The ITU/Zindi traces provide energy demand, PV and grid inputs, but
not observed diesel-order or delivery records. Tank inventory,
replenishment requests, pending deliveries and inventory transitions
are simulated under the stated logistics assumptions.
\end{revisedblock}

% ---------------------------------------------------------------
\section{Constraints}
\label{sec:constraints}
% ---------------------------------------------------------------

\begin{revisedblock}
Battery state of charge is bounded by
$0.20\leq\SoC_t\leq1.00$, with charge and discharge
power constrained by the site battery limits. The simulator
applies these limits through its fixed dispatch rule.

A DG ON request is feasible only if the available inventory is
at least $1.20\times2\,\mathrm{h}\times
P_{\mathrm{rated}}^{\DG}$ in kWh-equivalent units. This
is an activation reserve, not a continuous minimum ON-time
requirement.

A new order is not accepted while another order is pending.
The requested order is drawn from the two tank-relative sizes
and is subject to remaining tank headroom. Feasible action
masking excludes certain requests before selection; the
simulator also blocks or clips infeasible requests submitted
to its step interface. These rules define the feasible action
set $\mathcal C(s_t)$ used below.
\end{revisedblock}

% ---------------------------------------------------------------
\section{Objective Function}
\label{sec:objective}
% ---------------------------------------------------------------

\begin{revisedblock}
The implemented controller minimises the expected discounted
sum of the normalised step cost
\begin{equation}
J(\pi)=\mathbb E_{\pi}\!\left[
\sum_{t=0}^{T-1}\gamma^t c_t\right],
\label{eq:objective}
\end{equation}
where
\begin{equation}
c_t=
\frac{\alpha E_t^{\DG}+\beta E_t^{\Grid}+\lambda U_t}{\bar D}
+\gamma_r S_t+\mu V_t.
\label{eq:implemented_step_cost}
\end{equation}
Here $E_t^{\DG}$ is actual diesel-generated energy,
$E_t^{\Grid}$ is total grid energy, and $U_t$ is unmet
load energy. $S_t$ indicates a stockout coinciding with
unmet load after fuel use and before delivery, while $V_t$
indicates a physical-bound violation detected by the
simulator. The coefficients are $\alpha=1.0$, $\beta=0.4$,
$\lambda=100$, $\gamma_r=200$ and $\mu=20$.
The implemented discount factor is $\gamma=0.99$.

The cost is the negative of the reward used for PPO
training (Section~\ref{sec:reward_design}). The large
unmet-load coefficient prioritises service reliability.
A reward received 336 hourly steps in the future has
weight $0.99^{336}\approx0.034$ relative to an immediate
reward. Discounting alone does not establish the cause
of the observed replenishment behaviour.

\textbf{Economic extensions not optimised.}
An economic dispatch-and-inventory formulation could
add order-placement and ordered-volume costs, inventory
holding charges, battery degradation and carbon costs.
None of these is an explicit term in the implemented
reward; no claim of minimum monetary operating cost is
made. These extensions are discussed in
Chapter~\ref{chap:conclusion}.
\end{revisedblock}

% ---------------------------------------------------------------
\section{Formal CMDP Definition}
\label{sec:cmdp}
% ---------------------------------------------------------------

\begin{revisedblock}
The implemented control problem is represented as an episodic
feasibility-constrained Markov decision process,
\[
(\mathcal S,\mathcal A,\mathcal P,c,\mathcal C,\gamma).
\]
The agent receives the 11-scalar observation of
Section~\ref{sec:state}; its action set contains the six
generator--order combinations of Section~\ref{sec:action}.
The transition model includes site energy inputs,
fixed-priority dispatch and stochastic delivery.
The step cost $c_t$ is given in
Equation~\eqref{eq:implemented_step_cost}, the feasible
action set is $\mathcal C(s_t)$, and $\gamma=0.99$.
Future delivery times are not directly observed.

Training episodes use 720 hourly steps drawn from the
45-day training segment. Evaluation episodes use 360
hourly steps from the held-out 15-day segment; the
evaluation protocol is detailed in
Chapter~\ref{chap:methodology}.
\end{revisedblock}

% ---------------------------------------------------------------
\section{Unified versus Decomposed Architecture}
\label{sec:unified_vs_hier}
% ---------------------------------------------------------------

A natural modelling question is whether the diesel ordering decision
($Q_t$, multi-day timescale) and the energy dispatch decisions (generator
commitment and battery behaviour, hourly timescale) should be optimised
jointly within a single agent or decomposed into separate controllers
operating at their natural timescales. This question was raised explicitly
in the Phase-I reviewer comments and is directly addressed by the
experimental design of this work.

The unified formulation presented in this chapter---in which a single
joint agent controls diesel ordering and generator commitment while battery
behaviour is handled by the environment---is referred to as \textbf{RLInv}
throughout. The alternative \textbf{Track~B} architecture pairs a classical
$(s,S)$ inventory policy for diesel ordering with a separate PPO agent for
dispatch, allowing each sub-problem to be solved at its natural timescale.
The two architectures are evaluated empirically in
Chapter~\ref{chap:results}.

The comparison between unified and decomposed architectures allows the
following hypotheses to be tested.

\begin{revisedblock}
\textbf{H1 (Joint ordering).} Does a jointly learned
DG-and-ordering policy reduce EENS relative to a classical
$(s,S)$ ordering policy when dispatch configurations are
broadly comparable?

\textbf{H2 (Action masking).} Does feasibility-aware masking
improve reliability under stressed logistics, and does the
effect depend on how ordering is controlled?

\textbf{H3 (Logistics observations).} Does zeroing the five
inventory- and delivery-related observation channels reduce
reliability relative to the full observation? The A5 ablation
retains the 11-scalar observation shape and zeros those five
channels rather than deleting them.

These are testable questions, not assumptions that the
proposed controller must outperform its comparators.
The comparative results and their limitations are given in
Chapter~\ref{chap:results}.
\end{revisedblock}


% ---------------------------------------------------------------
\section{Research Objectives and Contributions}
\label{sec:objectives}
% ---------------------------------------------------------------

The formulation presented in this chapter serves three related objectives:
it enables joint dispatch and diesel-ordering control within a single
sequential decision framework; it incorporates stochastic delivery-pipeline
dynamics into the state representation, capturing replenishment uncertainty
that the energy-management formulations reviewed in Chapter 2 do not explicitly model; and it encodes
hard operational feasibility requirements directly in the decision problem,
rather than relying on penalty terms alone.

The primary contribution of this chapter is the formulation of telecom tower
energy management as a constrained sequential decision problem that jointly
models diesel inventory logistics and energy dispatch. The chapter defines the
unified state--action representation, the stochastic delivery-pipeline model,
and the experimental hypothesis structure used to evaluate the formulation.
The empirical evaluation of these hypotheses is presented in
Chapter~\ref{chap:results}.

\paragraph{Stakeholder impact.}
Rural telecom infrastructure connects remote populations to emergency
communications, mobile banking, and digital public services.  \revised{Telecom infrastructure in weak-grid areas depends on
reliable grid, battery and backup generation; energy and fuel
costs are material operating considerations
\citep{gsma2025b,trai2024}.}
Tower outages in these settings directly interrupt critical services and
disproportionately affect underserved communities.  Improving energy and
inventory management at these sites therefore has direct implications for
service reliability, operating cost, and carbon emissions.  The formulation
presented in this chapter targets precisely these drivers by jointly selecting DG commitment and replenishment actions under realistic supply-chain
uncertainty.